\subsection{PROPERTIES OF INTEGRALS}

The properties of the integrals of regulated functions are simply a translation, into the appropriate notation, of the properties of derivatives demonstrated in chap.~I.

In the first place, the formula (3) shows that no matter what the points \(x, y, z\) of \(I\) , one has

\[
\int_ {x} ^ {x} \mathbf {f} (t) d t = 0\tag{4}
\]

\[
\int_ {x} ^ {y} \mathbf {f} (t) d t + \int_ {y} ^ {x} \mathbf {f} (t) d t = 0\tag{5}
\]

\[
\int_ {x} ^ {y} \mathbf {f} (t) d t + \int_ {y} ^ {z} \mathbf {f} (t) d t + \int_ {z} ^ {x} \mathbf {f} (t) d t = 0\tag{6}
\]

From prop. 1 of II, p.~52 one has

\[
\int_ {x _ {0}} (\mathbf {f} + \mathbf {g}) = \int_ {x _ {0}} \mathbf {f} + \int_ {x _ {0}} \mathbf {g}\tag{7}
\]

and for every scalar \(k\)

\[
\int_ {x _ {0}} k \mathbf {f} = k \int_ {x _ {0}} \mathbf {f}.\tag{8}
\]

Let \(E, F\) be two complete normed spaces over \(\mathbf{R}\) , and \(\mathbf{u}\) a continuous linear map from \(E\) into \(F\) . If \(\mathbf{f}\) is a regulated function on \(I\) with values in \(E\) , then \(\mathbf{u} \circ \mathbf{f}\) is a regulated function on \(I\) with values in \(F\) (II, p.~6, cor. 2), and one has for \(a, b \in I\) (I, p.~13, prop. 2)

\[
\int_ {a} ^ {b} \mathbf {u} (\mathbf {f} (t)) d t = \mathbf {u} \left(\int_ {a} ^ {b} \mathbf {f} (t) d t\right).\tag{9}
\]

Now let E, F, G be three complete normed spaces over R, and \((\mathbf{x}, \mathbf{y}) \mapsto [\mathbf{x} \cdot \mathbf{y}]\) a continuous bilinear map of \(E \times F\) into G. Let f and g be two vector functions defined and continuous on I, taking their values in E and F respectively; suppose moreover that f and g are two primitives of regulated functions, which we denote by \(f'\) and \(g'\) by abuse of language (these functions are not actually guaranteed to be equal to the derivatives of f and g respectively except on the complement of a countable set). By prop. 3 of I, p.~6, the function \(\mathbf{h}(x) = [\mathbf{f}(x) \cdot \mathbf{g}(x)]\) has, at every point of the complement of a countable subset of I, a derivative equal to \([\mathbf{f}(x) \cdot \mathbf{g}'(x)] + [\mathbf{f}'(x) \cdot \mathbf{g}(x)]\) . Now, by the continuity of {[}x.y{]} and cor. 2 of II, p.~55, each of the functions \([f \cdot g']\) and \([f' \cdot g]\) is a regulated function on I; so one has the formula

\[
\int_ {a} ^ {b} [ \mathbf {f} ^ {\prime} (t). \mathbf {g} (t) ] d t = [ \mathbf {f} (t). \mathbf {g} (t) ] \big | _ {a} ^ {b} - \int_ {a} ^ {b} [ \mathbf {f} (t). \mathbf {g} ^ {\prime} (t) ] d t\tag{10}
\]

called the formula for integration by parts, which allows one to evaluate many primitives.

For example, the formula for integration by parts yields the following formula

\[
\int_ {x _ {0}} ^ {x} t \mathbf {f} ^ {\prime} (t) d t = t \mathbf {f} (t) \Big | _ {x _ {0}} ^ {x} - \int_ {x _ {0}} ^ {x} \mathbf {f} (t) d t
\]

so reducing the evaluation of primitives of one of the two functions \(\mathbf{f}(x)\) and \(x\mathbf{f}'(x)\) to the other.

Likewise, if \(\mathbf{f}\) and \(\mathbf{g}\) are \(n\) times differentiable on an interval I, and if \(\mathbf{f}^{(n)}\) and \(\mathbf{g}^{(n)}\) are regulated functions on I, then formula (5) of I, p.~21 is equivalent to the following:

\[
\begin{array}{l} \int_ {a} ^ {b} [ \mathbf {f} ^ {(n)} (t). \mathbf {g} (t) ] d t \\ = \left. \left(\sum_ {p = 0} ^ {n - 1} (- 1) ^ {p} [ \mathbf {f} ^ {(n - p - 1)} (t). \mathbf {g} ^ {(p)} (t) ]\right) \right| _ {a} ^ {b} + (- 1) ^ {n} \int_ {a} ^ {b} [ \mathbf {f} (t). \mathbf {g} ^ {(n)} (t) ] d t \end{array}\tag{11}
\]

which one calls the formula for integration by parts of order n.

Let us now translate the formula for differentiation of composite functions (I, p.~9, prop. 5). Let \(f\) be a real function defined and continuous on I, which is the primitive of a regulated function on I (which we again write as \(f'\) by abuse of language); let, moreover, \(\mathbf{g}\) be a continuous vector function (with values in a complete normed space) on an open interval J containing \(f(\mathrm{I})\) ; if \(\mathbf{h}\) denotes an arbitrary primitive of \(\mathbf{g}\) on J, then \(\mathbf{h}\) admits a derivative equal to \(\mathbf{g}\) at each point of J (II, p.~56, prop. 3); thus the composite function \(\mathbf{h} \circ f\) admits a derivative equal to \(\mathbf{g}(f(x))f'(x)\) at all the points of the complement (with respect to I) of a countable subset of I (I, p.~9, prop. 5); since the function \(\mathbf{g}(f(x))f'(x)\) is regulated (II, p.~55, cor. 2), one can write the formula

\[
\int_ {a} ^ {b} \mathbf {g} (f (t)) f ^ {\prime} (t) d t = \int_ {f (a)} ^ {f (b)} \mathbf {g} (u) d u\tag{12}
\]

called the formula for change of variables, which also facilitates the evaluation of primitives.

If, for example, one takes \(f(x) = x^{2}\) , one sees that the formula (12) reduces from one to the other the evaluation of the primitives of the functions \(\mathbf{g}(x)\) and \(x\mathbf{g}(x^{2})\) .

To translate the mean value theorem (I, p.~14, th. 1) for primitives of real regulated functions, we first remark that a real regulated function \(f\) on a compact interval I is bounded on I; let J be the set of points of I where \(f\) is continuous, and put \(m = \inf_{x \in \mathrm{J}} f(x)\) , \(\mathrm{M} = \sup_{x \in \mathrm{J}} f(x)\) ; one knows (II, p.~54, th. 3) that \(\mathrm{I} \cap \complement \mathrm{J}\) is countable; further, if B is the complement, with respect to I, of any countable subset of I, and \(m' = \inf_{x \in \mathrm{B}} f(x)\) , \(\mathrm{M}' = \sup_{x \in \mathrm{B}} f(x)\) , then one has \(m' \leqslant m \leqslant M \leqslant M'\) : indeed, for every point \(x \in J\) , there are points y of B arbitrarily close to x, whence \(m' \leqslant f(y) \leqslant M'\) ; since f is continuous at the point x one sees, making y approach x (y remaining in B) that \(m' \leqslant f(x) \leqslant M'\) , which proves our assertion. This being so, translating the mean value theorem gives the following proposition:

PROPOSITION 6 (theorem of the mean). Let \(f\) be a real regulated function on a compact interval \(\mathrm{I} = [a, b]\) ; if \(\mathrm{J}\) is the set of points of \(\mathrm{I}\) where \(f\) is continuous, and \(m = \inf_{x \in \mathrm{J}} f(x)\) , \(\mathrm{M} = \sup_{x \in \mathrm{J}} f(x)\) , then

\[
m <   \frac {1}{b - a} \int_ {a} ^ {b} \mathbf {f} (t) d t <   \mathbf {M}\tag{13}
\]

except when \(f\) is constant on \(\mathbf{J}\) , in which case the three members of (13) are equal.

In other words, the mean of the regulated function \(f\) in I lies between the bounds of \(f\) over the subset of I where \(f\) is continuous.

COROLLARY 1. If a real regulated function \(f\) on \(I\) is such that \(f(x) \geqslant 0\) at the points where \(f\) is continuous, then \(\frac{1}{b - a} \int_{a}^{b} f(t) dt > 0\) unless \(f(x) = 0\) at the points where \(f\) is continuous.

COROLLARY 2. Let f and g be two real regulated functions on I, such that \(g(x) \geqslant 0\) at the points where g is continuous; if m and M are the greatest lower bound and least upper bound of f over the set of points where f is continuous, then

\[
\frac {m}{b - a} \int_ {a} ^ {b} g (t) d t \leqslant \frac {1}{b - a} \int_ {a} ^ {b} f (t) g (t) d t \leqslant \frac {\mathrm{M}}{b - a} \int_ {a} ^ {b} g (t) d t.\tag{14}
\]

The first two terms (resp. the two last) are unequal unless \(g(x)(f(x) - m) = 0\) (resp. \(g(x)(f(x) - \mathbf{M}) = 0\) ) at every point where \(f\) and \(g\) are continuous.

For vector functions the mean value theorem (I, p.~15, th. 2) yields the following proposition:

PROPOSITION 7. Let \(\mathbf{f}\) be a regulated vector function on a compact interval \(\mathrm{I} = [a, b]\) , with values in a complete normed space \(\mathrm{E}\) , and let \(g\) be a real regulated function on \(\mathrm{I}\) , such that \(g(x) \geqslant 0\) at the points where \(g\) is continuous; in these circumstances

\[
\left\| \int_ {a} ^ {b} \mathbf {f} (t) g (t) d t \right\| \leqslant \int_ {a} ^ {b} \| \mathbf {f} (t) \| g (t) d t.\tag{15}
\]

In particular,

\[
\left\| \int_ {a} ^ {b} \mathbf {f} (t) d t \right\| \leqslant \int_ {a} ^ {b} \| \mathbf {f} (t) \| d t.\tag{16}
\]

